% =============================================================================
%  Chapter 2 — Cơ sở lý thuyết & Các công trình liên quan
%  Section: Mở đầu chương
% =============================================================================

Chương này trình bày các cơ sở lý thuyết nền tảng phục vụ cho đề tài, bao gồm các phương pháp biểu diễn tri thức có cấu trúc như Ontology và Knowledge Graph, cùng với các kiến trúc truy hồi – tạo sinh hiện đại như Retrieval-Augmented Generation (RAG) và GraphRAG. Các khái niệm này đóng vai trò quan trọng trong việc tổ chức, khai thác và suy luận tri thức, đồng thời là nền tảng cho hệ thống được đề xuất trong khóa luận. Bên cạnh đó, chương cũng tổng hợp và phân tích các công trình nghiên cứu liên quan nhằm làm rõ bối cảnh và định hướng nghiên cứu.
